\section{Related Work}

\paragraph{Latent world models and planning.}
Ha and Schmidhuber~\citep{ha2018world} established the template of learning a
compressed latent dynamics model and acting through it. PlaNet~\citep{hafner2019planet}
planned directly in a learned latent space, the Dreamer line~\citep{hafner2023dreamerv3}
scaled latent imagination to diverse domains, and MuZero~\citep{schrittwieser2020muzero}
showed that a learned model is sufficient for strong planning. These models
treat the action as an input to a monolithic transition. AC-LSCM instead
factorizes the transition per latent factor over a learned sparse graph and
treats actions as interventions that replace individual mechanisms.

\paragraph{Causal structure in dynamics models.}
Causal Dynamics Learning~\citep{wang2022cdl} learns a sparse causal graph
between state variables and actions in order to obtain task-independent state
abstractions, and Ke et al.~\citep{ke2021systematic} systematically evaluate
causal-discovery-based world models in model-based reinforcement learning.
The simplified AC-LSCM is close in spirit to these works: a sparse, lagged
graph over state factors learned jointly with per-factor mechanisms.

\paragraph{Causal representation learning.}
Sch\"olkopf et al.~\citep{scholkopf2021causal} survey the program of learning
representations aligned with causal mechanisms. CITRIS~\citep{lippe2022citris}
proves identifiability of causal factors from temporal sequences with known
intervention targets, which is the regime of our experiments; mechanism
sparsity~\citep{lachapelle2022disentanglement} and weak supervision from
paired interventions~\citep{brehmer2022weakly} provide related guarantees.
We do not contribute identifiability results; our encoder is a plain MLP and
our latents are anchored by the known action targets.

\paragraph{Continuous DAG learning.}
NOTEARS~\citep{zheng2018notears} reformulated DAG learning as continuous
optimization with a smooth acyclicity function, and DAGMA~\citep{bello2022dagma}
improved its numerical behaviour with a log-determinant characterization.
Our negative result is specific to how these penalties were combined with a
relaxed adjacency and a sparsity term in a lagged model: there, the
acyclicity term is unnecessary and, with the original weights, the combined
penalty removed every edge.

\paragraph{Safe action selection.}
Shielding~\citep{alshiekh2018shielding} blocks actions whose consequences
would violate a formal safety specification, using a model of the environment.
Algorithm~\ref{alg:planning} is a learned, approximate analogue: the shield is
the safety predicate evaluated on the forward model's simulated trajectory,
with deferral when no action passes.

\paragraph{World models and risk for LLM agents.}
RAP~\citep{hao2023rap} uses an LLM as its own world model for planning, and
WebDreamer~\citep{gu2024webdreamer} uses an LLM to simulate the outcome of
candidate web actions before executing them. ToolEmu~\citep{ruan2024toolemu}
uses an LM-emulated sandbox to identify risky agent actions. These systems
obtain the forward model from the language model itself; our question is what
structure a forward model needs in order to get the consequences of
interventions right, studied in a setting where ground truth is available.
CLadder~\citep{jin2023cladder} and CounterBench~\citep{chen2025counterbench}
document weaknesses of current LLMs on formal causal queries and motivate the
problem, but we have not evaluated on them.

\section{Limitations and Open Questions}
\label{sec:limitations}

\paragraph{Synthetic, low-dimensional setting.}
All results are on synthetic SCMs with $K \le 20$ scalar factors, observations
that are a noisy copy of the latent state, single-factor actions with known
targets, and small MLPs. The ground-truth mechanisms are additive-noise and
lagged, which is exactly the model class of the simplified AC-LSCM. Whether
the gains survive richer observations, unknown intervention targets,
instantaneous effects, or misspecified mechanisms is untested.

\paragraph{No language-domain evaluation.}
We have not connected AC-LSCM to a language model or evaluated it on
text-based agent tasks or causal-reasoning benchmarks. The paper's motivation
concerns LLM agents; its evidence concerns forward models in synthetic
environments. Bridging that gap is the most important open step.

\paragraph{Counterfactual supervision.}
The counterfactual loss requires ground-truth counterfactual latents, which
are available only in simulation and give AC-LSCM information about the true
state coordinates that the baselines do not receive. Without it, AC-LSCM's
consequence-task safety is on par with the best baseline
(Sections~\ref{sec:consequence} and~\ref{sec:ablations}).

\paragraph{Task design.}
The consequence task uses one hand-chosen safety factor and one goal factor
on one graph, a fixed horizon of three steps, and candidate sets constructed
so that a fixed fraction of episodes has no safe action. Candidates whose
realized outcome falls within a margin of a threshold are excluded so that
labels are decided by the action rather than by noise. Other choices would
give different absolute numbers.

\paragraph{Evaluation protocol.}
Early stopping uses a validation split taken from the training trajectory.
The simplified model's settings came from the first-version ablations, one
short pilot run, and one further change (the straight-through mask) chosen on
prediction error after its first full runs; no setting was tuned on the agent
tasks, and the baselines keep the hyperparameters of the first version. The
agent-task configurations themselves (thresholds, horizon, margin, candidate
counts) were fixed once using only the oracle and two model-free reference
predictors.

\section{Conclusion}

Agents that act on real systems should predict what their actions will cause
before acting. We studied a structured forward model for this purpose and,
re-examining an earlier version of this work, found that its original
specification learned no causal structure and that its headline agent result
measured extrapolation of action values rather than modeling of consequences.
With those problems fixed, the evidence is narrower and, we believe, more
useful. A factorized, intervention-aware forward model gives the best
counterfactual predictions in our synthetic setting. In planning, what
matters is whether the model captures the indirect paths that lead to unsafe
states: a model that keeps all true edges matches an oracle, a model that
prunes one weak edge into the safety factor loses much of that advantage, and
the explicit do-operator contributes nothing measurable. Without ground-truth
counterfactual supervision, the architecture is on par with a well-tuned
unstructured latent model. For safety-critical use this argues for
conservative (denser) structure and for evaluation protocols that test
downstream consequences, in range of the training data, with every model
judged in the same space. Connecting such forward models to language-model
agents remains future work.


\section*{Code and Reproducibility}
Code, configurations, per-seed result files, and the scripts that produce
every table are available at
\url{https://github.com/mplusm/AC-LSCM-Outcome-Aware-Research-Paper}.
The first-version experiments ran on a single NVIDIA Tesla T4 in fp32
(under 50 GPU-hours). The experiments added in this version ran on a
4-core CPU. All runs use deterministic seeding.

\section*{Acknowledgments}
The author thanks the open research community working on causal
representation learning, world models, and agent safety, whose published
work and benchmarks made this study possible.

\begin{thebibliography}{99}

\bibitem[Alshiekh et~al.(2018)]{alshiekh2018shielding}
M.~Alshiekh, R.~Bloem, R.~Ehlers, B.~K{\"o}nighofer, S.~Niekum, and U.~Topcu.
\newblock Safe reinforcement learning via shielding.
\newblock In \emph{Proceedings of the AAAI Conference on Artificial Intelligence}, 2018.
\newblock arXiv:1708.08611.

\bibitem[Bello et~al.(2022)]{bello2022dagma}
K.~Bello, B.~Aragam, and P.~Ravikumar.
\newblock {DAGMA}: Learning {DAG}s via {M}-matrices and a log-determinant acyclicity characterization.
\newblock In \emph{Advances in Neural Information Processing Systems}, volume~35, 2022.

\bibitem[Boutilier et~al.(2000)]{boutilier2000factored}
C.~Boutilier, R.~Dearden, and M.~Goldszmidt.
\newblock Stochastic dynamic programming with factored representations.
\newblock \emph{Artificial Intelligence}, 121(1--2):49--107, 2000.

\bibitem[Brehmer et~al.(2022)]{brehmer2022weakly}
J.~Brehmer, P.~de~Haan, P.~Lippe, and T.~Cohen.
\newblock Weakly supervised causal representation learning.
\newblock In \emph{Advances in Neural Information Processing Systems}, volume~35, 2022.
\newblock arXiv:2203.16437.

\bibitem[Chen et~al.(2025)]{chen2025counterbench}
Y.~Chen, V.~K. Singh, J.~Ma, and R.~Tang.
\newblock {CounterBench}: Evaluating and improving counterfactual reasoning in large language models.
\newblock arXiv:2502.11008, 2025.

\bibitem[Gu et~al.(2024)]{gu2024webdreamer}
Y.~Gu, K.~Zhang, Y.~Ning, B.~Zheng, B.~Gou, T.~Xue, C.~Chang, S.~Srivastava, Y.~Xie, P.~Qi, H.~Sun, and Y.~Su.
\newblock Is your {LLM} secretly a world model of the internet? {M}odel-based planning for web agents.
\newblock arXiv:2411.06559, 2024.

\bibitem[Ha and Schmidhuber(2018)]{ha2018world}
D.~Ha and J.~Schmidhuber.
\newblock World models.
\newblock arXiv:1803.10122, 2018.

\bibitem[Hafner et~al.(2019)]{hafner2019planet}
D.~Hafner, T.~Lillicrap, I.~Fischer, R.~Villegas, D.~Ha, H.~Lee, and J.~Davidson.
\newblock Learning latent dynamics for planning from pixels.
\newblock In \emph{International Conference on Machine Learning}, 2019.
\newblock arXiv:1811.04551.

\bibitem[Hafner et~al.(2023)]{hafner2023dreamerv3}
D.~Hafner, J.~Pasukonis, J.~Ba, and T.~Lillicrap.
\newblock Mastering diverse domains through world models.
\newblock arXiv:2301.04104, 2023.

\bibitem[Hao et~al.(2023)]{hao2023rap}
S.~Hao, Y.~Gu, H.~Ma, J.~J. Hong, Z.~Wang, D.~Z. Wang, and Z.~Hu.
\newblock Reasoning with language model is planning with world model.
\newblock In \emph{Proceedings of the Conference on Empirical Methods in Natural Language Processing}, 2023.
\newblock arXiv:2305.14992.

\bibitem[Jin et~al.(2023)]{jin2023cladder}
Z.~Jin, Y.~Chen, F.~Leeb, L.~Gresele, O.~Kamal, Z.~Lyu, K.~Blin, F.~Gonzalez~Adauto, M.~Kleiman-Weiner, M.~Sachan, and B.~Sch{\"o}lkopf.
\newblock {CLadder}: Assessing causal reasoning in language models.
\newblock In \emph{Advances in Neural Information Processing Systems}, volume~36, 2023.

\bibitem[Ke et~al.(2021)]{ke2021systematic}
N.~R. Ke, A.~Didolkar, S.~Mittal, A.~Goyal, G.~Lajoie, S.~Bauer, D.~Rezende, Y.~Bengio, M.~Mozer, and C.~Pal.
\newblock Systematic evaluation of causal discovery in visual model based reinforcement learning.
\newblock arXiv:2107.00848, 2021.

\bibitem[Lachapelle et~al.(2022)]{lachapelle2022disentanglement}
S.~Lachapelle, P.~Rodr{\'i}guez~L{\'o}pez, Y.~Sharma, K.~Everett, R.~Le~Priol, A.~Lacoste, and S.~Lacoste-Julien.
\newblock Disentanglement via mechanism sparsity regularization: A new principle for nonlinear {ICA}.
\newblock In \emph{Conference on Causal Learning and Reasoning}, 2022.
\newblock arXiv:2107.10098.

\bibitem[Lippe et~al.(2022)]{lippe2022citris}
P.~Lippe, S.~Magliacane, S.~L{\"o}we, Y.~M. Asano, T.~Cohen, and E.~Gavves.
\newblock {CITRIS}: Causal identifiability from temporal intervened sequences.
\newblock In \emph{International Conference on Machine Learning}, 2022.
\newblock arXiv:2202.03169.

\bibitem[Madapathi(2026)]{madapathi2026zenodo}
M.~Madapathi.
\newblock Outcome-aware agents: From token prediction to action consequence modeling (first version).
\newblock Zenodo, 2026.
\newblock \url{https://doi.org/10.5281/zenodo.20379456}.

\bibitem[Pearl(2009)]{pearl2009causality}
J.~Pearl.
\newblock \emph{Causality: Models, Reasoning, and Inference}.
\newblock Cambridge University Press, 2nd edition, 2009.

\bibitem[Pitis et~al.(2020)]{pitis2020coda}
S.~Pitis, E.~Creager, and A.~Garg.
\newblock Counterfactual data augmentation using locally factored dynamics.
\newblock In \emph{Advances in Neural Information Processing Systems}, volume~33, 2020.
\newblock arXiv:2007.02863.

\bibitem[Ruan et~al.(2024)]{ruan2024toolemu}
Y.~Ruan, H.~Dong, A.~Wang, S.~Pitis, Y.~Zhou, J.~Ba, Y.~Dubois, C.~J. Maddison, and T.~Hashimoto.
\newblock Identifying the risks of {LM} agents with an {LM}-emulated sandbox.
\newblock In \emph{International Conference on Learning Representations}, 2024.
\newblock arXiv:2309.15817.

\bibitem[Sch{\"o}lkopf et~al.(2021)]{scholkopf2021causal}
B.~Sch{\"o}lkopf, F.~Locatello, S.~Bauer, N.~R. Ke, N.~Kalchbrenner, A.~Goyal, and Y.~Bengio.
\newblock Toward causal representation learning.
\newblock \emph{Proceedings of the IEEE}, 109(5):612--634, 2021.

\bibitem[Schrittwieser et~al.(2020)]{schrittwieser2020muzero}
J.~Schrittwieser, I.~Antonoglou, T.~Hubert, K.~Simonyan, L.~Sifre, S.~Schmitt, A.~Guez, E.~Lockhart, D.~Hassabis, T.~Graepel, T.~Lillicrap, and D.~Silver.
\newblock Mastering {A}tari, {G}o, chess and shogi by planning with a learned model.
\newblock \emph{Nature}, 588:604--609, 2020.

\bibitem[Wang et~al.(2022)]{wang2022cdl}
Z.~Wang, X.~Xiao, Z.~Xu, Y.~Zhu, and P.~Stone.
\newblock Causal dynamics learning for task-independent state abstraction.
\newblock In \emph{International Conference on Machine Learning}, 2022.
\newblock arXiv:2206.13452.

\bibitem[Zheng et~al.(2018)]{zheng2018notears}
X.~Zheng, B.~Aragam, P.~Ravikumar, and E.~P. Xing.
\newblock {DAG}s with {NO TEARS}: Continuous optimization for structure learning.
\newblock In \emph{Advances in Neural Information Processing Systems}, volume~31, 2018.

\end{thebibliography}


\appendix

\section{Experimental Details}
\label{app:details}

\paragraph{Data.}
Each environment is an SCM over $K$ scalar factors with lagged additive-noise
mechanisms $z^i_{t+1} = \tanh\big(\sum_j W_{ij} z^j_t\big) + \epsilon^i_t$,
$\epsilon^i_t \sim \mathcal{N}(0, 0.1^2)$, with edge weights drawn uniformly
from $\pm[0.5, 2]$. Observations are $\x_t = \z_t + \mathcal{N}(0, 0.1^2)$.
Graphs: chain, fork, and collider with $K=10$, and Erd\H{o}s--R\'enyi graphs
with edge probability $0.3$ for $K=10$ (8 edges) and $K=20$ (56 edges). Each
step applies, with probability $0.5$, one intervention $\doop(z^i = v)$ with
$i$ uniform and $v \sim U(-1.5, 1.5)$. Each training set is one trajectory of
$N = 10{,}000$ steps; each test set is a separate trajectory of 500 steps.
Counterfactual targets reuse the same noise with a different random
single-factor action.

\paragraph{Models.}
All latent models share an encoder and decoder (two-layer MLPs, hidden size
128) with a $K$-dimensional latent. AC-LSCM uses one transition MLP per factor
(input $K{+}1$, hidden size 32) and a Gumbel-sigmoid adjacency with
temperature $0.5$. Vanilla VAE uses a single MLP transition (hidden 128) on the
concatenation of latent and action; Dreamer-style uses a GRU cell with hidden
size $K$. The Transformer embeds each of the $2K$ scalars of $(\x_t, a_t)$ as a
token ($d = 128$, 2 pre-norm layers, 4 heads, feed-forward 512), mean-pools,
and predicts $\x_{t+1}$ with an MLP head. Missing action components are
encoded as $0$ for all baselines.

\paragraph{Optimization.}
Adam with learning rate $10^{-3}$, batch size 64, gradient clipping at 1.0,
up to 200 epochs with early stopping (patience 20) on the validation loss.
Loss weights $\beta = (1, 1, 0.5, 1)$; the graph-loss weight is ramped
linearly from 0 over the first 50 epochs. Original AC-LSCM: $\lambda_1 = 0.1$,
$\lambda_2 = 1$ (NOTEARS), adjacency logits initialized at 0, hinge margin
$\gamma = 1$. Simplified AC-LSCM: $\lambda_1 = 0.01$, $\lambda_2 = 0$, logits
initialized at $+2$, no hinge, exact abduction, straight-through hard
Gumbel-sigmoid adjacency samples.

\paragraph{Compute.}
The experiments of this version (61 training runs) ran on a shared 4-core
CPU without a GPU; a run of AC-LSCM on ER-10 takes roughly 15 to 45 minutes
there, and a run of the Transformer roughly 90 minutes.

\paragraph{Metrics.}
Intervention MSE is the observation-space error of the one-step prediction
under the logged action on the test trajectory. Counterfactual MSE is the
error of the counterfactual prediction against the ground-truth
counterfactual observation. Structural Hamming distance (SHD) counts
mismatched entries between the thresholded learned adjacency and the true
one. All tables report the mean and sample standard deviation over seeds.

\section{Relation to the First Version}
\label{app:v1}

The first version~\citep{madapathi2026zenodo} reported, on the direct-effect
task with predicates evaluated in latent space, a safety-violation rate of
$0.005 \pm 0.019$ for AC-LSCM over 13 seeds (8 trained on $N{=}10{,}000$
samples and 5 on $N{=}2{,}000$, pooled) against $0.180 \pm 0.022$ for the
Transformer over 3 seeds, and attributed the gain to the do-operator and
abduction. This version re-runs that task with the same code and finds:
(i) the original model's graph is empty in every run, so its predictions do
not depend on the state; (ii) the four ``failed'' seeds of the first version
(three ``passive'' and one ``mis-ranked'') are an artifact of thresholding
latents that are not aligned with observation coordinates and disappear when
predictions are decoded; (iii) the Transformer's violations come from
extrapolating action values outside the training range; and (iv) the
do-operator ablation has no effect. The ablation conclusions of the first
version that the hinge term and the NOTEARS term should be removed are
confirmed; the claims about the do-operator, about precise structure
recovery (SHD $8$), and about the $36\times$ safety improvement are withdrawn.
The first-version per-seed results remain in the repository under
\texttt{results/}; this version's are under \texttt{results\_v2/}.


\end{document}
